\subsection{平面三角学。}

在这一小节中，我们选取代数 A(Φ) 的一个生成元 w，使得 \(w^{2} \in A\) 。这样的生成元确定到相差一个位似（第 1 小节，注 2），因此 A 的元素 \(w^{2}\) （我们把它记作 δ）确定到模 A 的非零元素的平方所成的乘法子群 (A*)²。

注. — 当 -1 属于所说的模 (A*)² 的类时，通常选取 w 使 \(w^{2} = -1\) ，这就把它确定到相差一个符号。当这个类含 1 而不含 -1 时，通常选取 w 使 \(w^{2} = 1\) ，这同样把它确定到相差一个符号。

这样，每个元素 \(\nu\) （\(S^{+}\) 的）都可以用一种且仅一种方式写成

\[
\nu = c _ {w} (\nu) + s _ {w} (\nu) \varpi\tag{1}
\]

其中 \(c_w(v)\) 、\(s_w(v)\) 属于 A；元素 \(s_w(v)/c_w(v)\) （属于射影域 \(\tilde{\mathbf{A}}\) ；第二章，2 \(^{e}\) 版，附录 III，第 5 小节）记作 \(t_w(v)\) ；它只依赖于 \(v\) 模 H 的类；于是 \(t_w\) 通过取商定义了一个从 S \(^{+}\) /H 到射影域 \(\tilde{\mathbf{A}}\) 中的映射，约定仍把它记作 \(t_w\) （滥用记号）。我们常以 \(c\) 、\(s\) 与 \(t\) 代替 \(c_w\) 、\(s_w\) 与 \(t_w\) 。我们有 \(c_{-w} = c_w\) ，\(s_{-w} = -s_w\) 与 \(t_{-w} = -t_w\) 。

命题 4. — a) 当 \(w^2 = \delta\) 不是 A 中的平方时（即当 E 不含迷向直线时），从 S⁺/H 到 \(\tilde{\mathbf{A}}\) 的映射 t 是双射。

\begin{enumerate}
\def\labelenumi{\alph{enumi})}
\setcounter{enumi}{1}
\item
  当 δ 是 A 中元素 γ 的平方时，映射 t 是从 S⁺/H 到 \(\tilde{\mathbf{A}}\) 除去元素 1/γ 与 -1/γ 后所成的集合上的双射。
\item
  把 \(\mathrm{S}^{+}/\mathrm{H}\) 写成加法，则对 \(\varphi, \varphi'\) （\(\mathrm{S}^{+}/\mathrm{H}\) 中的），有
\end{enumerate}

\[
t (\varphi + \varphi^ {\prime}) = (t (\varphi) + t (\varphi^ {\prime})) / (1 + \delta t (\varphi) t (\varphi^ {\prime}))\tag{2}
\]

当 \(t(\varphi)\) 与 \(t(\varphi')\) 有限且 \(1 + t(\varphi)t(\varphi')\) \(\neq 0\) 时。

事实上，由于 \(S^{+}/H\) 是把 \(A(\Phi)\) 看作 A 上向量平面时去掉 0 的那些直线所成的集合，t 是单射。另一方面，要使元素 \(a + bw\) （ \(a \in A\) ，\(b \in A\) ，属于 \(A(\Phi)\) ）是一个正相似变换，必须且只需它可逆，即 \(N(a + bw) = a^{2} - \delta b^{2} \neq 0\) ，亦即 \((b/a)^{2} \neq 1/\delta\) ；这就证明了 a) 与 b) 中的满射性断言。最后，相似变换 \(1 + t(\varphi)w\) 与 \(1 + t(\varphi')w\) 的乘积是相似变换 \(1 + \delta t(\varphi)t(\varphi') + (t(\varphi) + t(\varphi'))w\) ，这就证明了 c)。

命题 5. --- 把 O \(^{+}\) 写成加法。对 O \(^{+}\) 的每对元素 θ、θ'，我们有

\begin{enumerate}
\def\labelenumi{(\arabic{enumi})}
\setcounter{enumi}{2}
\tightlist
\item
\end{enumerate}

\[
c (\theta) ^ {2} - \delta s (\theta) ^ {2} = 1\tag{4}
\]

\[
c (\theta + \theta^ {\prime}) = c (\theta) c (\theta^ {\prime}) + \delta s (\theta) s (\theta^ {\prime})\tag{5}
\]

\[
s (\theta + \theta^ {\prime}) = s (\theta) c (\theta^ {\prime}) + c (\theta) s (\theta^ {\prime}).
\]

关系 (3) 确实表达了 \(\mathrm{N}(c(\theta)+s(\theta)\omega)=1\) 。要证明 (4) 与 (5)，只需在 \(\mathrm{A}(\Phi)\) 中计算旋转 \(c(\theta)+s(\theta)\omega\) 与 \(c(\theta')+s(\theta')\omega\) 的乘积。

命题 6. — 设 d 是命题 3 中定义的从 S+/H 到 O+ 上的同构。对 S+/H 的每个使 t = t(φ) 为有限的元素 φ，有

\[
s (d (\varphi)) = 2 t / (1 - \delta t ^ {2}), \quad c (d (\varphi)) = (1 + \delta t ^ {2}) / (1 - \delta t ^ {2}).\tag{6}
\]

事实上，\(\varphi\) 是相似变换 \(1 + tw\) 模 H 的类，于是旋转 \(d(\varphi)\) 就是 \((1 + tw)^2/\mathrm{N}(1 + tw) = (1 + \delta t^2 + 2tw)/(1 - \delta t^2)\) （命题 3），这就证明了 (6)。

推论. --- 对 O⁺ 的每个使 \(t(\theta)\) 为有限的元素 θ，我们有

\[
s (2 \theta) = 2 t (\theta) / (1 - \delta t (\theta) ^ {2}), \quad c (2 \theta) = (1 + \delta t (\theta) ^ {2}) / (1 - \delta t (\theta) ^ {2}). \tag {7}
\]

对每个元素 \(\varphi\) （\(S^{+}/H\) 中的），只要 \(t(\varphi)\) 有限且 \(1 + \delta t(\varphi)^{2} \neq 0\) ，就有

\[
t (2 \varphi) = 2 t (\varphi) / (1 + \delta t (\varphi) ^ {2}).\tag{8}
\]

事实上，这由命题 6 及命题 3 的推论立即得出。

注. --- 公式 (6) 对 \(t = \infty\) 仍然成立，只要把出现在右端的那些有理函数换成它们到射影域的典范开拓，这里射影域指 \(\tilde{\mathbf{A}}\) （第二章，2 \(^{\text{e}}\) 版，附录 III，第 5 小节）；事实上，若 \(t = \infty\) ，则 \(\varphi\) 是 \(w\) 的类，且 \(d(\varphi) = -1\) ，\(s(d(\varphi)) = 0\) ，\(c(d(\varphi)) = -1\) ；而这些正是右端各式在 \(t = \infty\) 时由其典范开拓取到的值。当 \(t(\theta) = \infty\) 时 (7) 同样成立，当 \(t(\varphi) = \infty\) 或 \(1 + \delta t(\varphi)^{2} = 0\) 时 (8) 亦然。类似地，当 \(t(\varphi)\) 与 \(t(\varphi')\) 中只有一个（例如 \(t(\varphi)\) ）为无穷时，公式 (2) 仍成立，只要把它的右端看作单独 \(t(\varphi)\) 的有理函数：事实上，相似变换 \(1 + t(\varphi')w\) 与 \(w\) 的乘积是 \(\delta t(\varphi') + w\) ，而 (2) 的右端的典范开拓所取的值是 \(1/\delta t(\varphi')\) 。最后，当 \(t(\varphi)\) 与 \(t(\varphi')\) 都有限且有 \(1 + \delta t(\varphi)t(\varphi') = 0\) 时，有 \(t(\varphi) + t(\varphi') = 0\) （否则 \(t(\varphi)^{2}\) 将等于 \(1/\delta\) ，这是不可能的（命题 4））；因此可以约定 (2) 的右端的值为 \(\infty\) ，而这个值确实就是左端的值。当 \(t(\varphi)\) 与 \(t(\varphi')\) 都是无穷时，(2) 的右端没有定义。

